\chapter{Conclusion}
\label{ch:conclusion}

The European Union spent two decades pricing carbon while exempting the
industries most likely to respond to the price. The definitive \acrshort{cbam}
regime, in force since January 2026, is the first serious attempt to close
that gap: the border levy replaces \gls{freeallocation} rather than
supplementing it, and for the first time the marginal incentive facing an
emissions-intensive European producer is not neutralised by a free allowance.

This thesis treats that switch as a natural experiment. Because coverage is
enumerated by \acrshort{cn} code, the treated set is defined by the regulation
rather than by the researcher, and near-neighbour codes just outside the annex
supply a control group at the same level of aggregation. Applying a
\acrshort{did} design to monthly Eurostat import records, the analysis asks
whether covered imports fell, whether sourcing shifted towards cleaner origins,
and who paid.

[Findings to be summarised once \Cref{ch:results} is complete.]

Two things are worth stating regardless of what the estimates show. A null
result would be informative rather than empty: the earlier \acrshort{ets}
literature could always attribute its nulls to free allocation, and that
explanation is no longer available. And the design outlives this paper --- each
future extension of Annex~I supplies another dated boundary, so the question
can be re-asked with more power as the regime matures.
